\chapter{Sinus Pain}

\section*{Definition}
Sinus pain (rhinosinusitis) is facial pain or pressure localized over the paranasal sinuses (maxillary, frontal, ethmoid, sphenoid) typically accompanied by nasal congestion, purulent discharge, and hyposmia lasting more than 10 days (acute) or 12 weeks (chronic).

\section*{Pathophysiology}
Ostial obstruction (narrowing of sinus drainage pathways) from mucosal edema leads to mucus accumulation, decreased oxygen tension, and impaired mucociliary clearance within the sinus cavity. Bacterial superinfection of trapped secretions triggers neutrophilic inflammation, increasing intrasinus pressure and producing characteristic pain.

\section*{Etiology}
\begin{itemize}
  \item \textbf{Acute viral:} Rhinovirus, influenza, adenovirus (90\% of acute sinusitis)
  \item \textbf{Acute bacterial:} Streptococcus pneumoniae, Haemophilus influenzae, Moraxella catarrhalis
  \item \textbf{Chronic:} Persistent inflammation with biofilm formation, likely polymicrobial
  \item \textbf{Fungal:} Allergic fungal rhinosinusitis, acute invasive fungal sinusitis (immunocompromised)
  \item \textbf{Odontogenic:} Maxillary sinusitis from periapical dental abscess or molar root infection
  \item \textbf{Anatomic:} Septal deviation, concha bullosa, nasal polyps obstructing sinus ostia
\end{itemize}

\section*{Indications for Evaluation}
Seek care if:
\begin{itemize}
  \item Facial pain with purulent nasal discharge lasting > 10 days or worsening after 5--7 days (double sickening)
  \item Severe unilateral facial pain or dental pain (suspect odontogenic source)
  \item Periorbital edema, proptosis, or visual changes (suspect orbital cellulitis)
  \item Headache with neurologic signs (meningitis or intracranial extension)
  \item Immunocompromised patient with fever and facial pain (invasive fungal sinusitis)
\end{itemize}

\section*{Related Symptoms}
\begin{itemize}
  \item Facial pain/pressure (maxillary: cheek and teeth; frontal: forehead; ethmoid: medial canthus; sphenoid: vertex or occiput)
  \item Nasal congestion and obstruction
  \item Purulent anterior or posterior rhinorrhea
  \item Hyposmia or anosmia
  \item Fever, fatigue, and malaise
\end{itemize}
\textbf{Note:} Isolated facial pain without nasal congestion, discharge, or endoscopic findings is more likely due to tension headache, migraine, or temporomandibular dysfunction than sinusitis.

\section*{Self-Care / Home Management}
\begin{itemize}
  \item Saline nasal irrigation (Neti pot) 2--4 times daily
  \item Warm compresses over the affected sinuses
  \item Steam inhalation for 10--15 minutes (bowl of hot water with towel)
  \item Over-the-counter analgesics (acetaminophen, ibuprofen) for pain
  \item Intranasal corticosteroid sprays (fluticasone, mometasone) to reduce inflammation
  \item Topical decongestant sprays (oxymetazoline) limited to 3 days to avoid rhinitis medicamentosa
  \item Oral decongestants (pseudoephedrine 60 mg q4--6h) if no contraindication (hypertension, glaucoma)
\end{itemize}

\section*{Diagnostic Approach}
\begin{itemize}
  \item Nasal endoscopy to visualize purulence from sinus ostia, mucosal edema, and polyps
  \item Sinus CT without contrast (gold standard) for confirming mucosal thickening, air-fluid levels, and bony changes
  \item Sinus MRI for suspected fungal sinusitis or neoplasm
  \item Cultures (endoscopic-guided swab or sinus aspirate) for refractory acute sinusitis
  \item Allergy testing if chronic sinusitis with allergic component
  \item Dental evaluation with panoramic radiograph if odontogenic source suspected
\end{itemize}

\section*{Treatment / Management Overview}
\begin{itemize}
  \item \textbf{Acute viral:} Symptomatic management only; antibiotics not indicated
  \item \textbf{Acute bacterial:} Amoxicillin-clavulanate 875/125 mg BID for 5--7 days (first-line); doxycycline or respiratory fluoroquinolone for penicillin allergy
  \item \textbf{Chronic:} Intranasal corticosteroids, prolonged antibiotics (3--4 weeks if culture-directed), saline irrigations
  \item \textbf{Odontogenic:} Dental extraction or root canal plus anaerobic coverage (metronidazole or clindamycin)
  \item \textbf{Functional endoscopic sinus surgery (FESS):} For refractory chronic sinusitis, mucocele, or antrochoanal polyp
  \item \textbf{Balloon sinuplasty:} Minimally invasive dilation of sinus ostia (selected chronic cases)
  \item \textbf{Biologics:} Dupilumab for severe CRSwNP
\end{itemize}

\clearpage
